\section{Background and Motivation}
Modern discrete-manufacturing facilities increasingly rely on dense sensor deployments to
monitor equipment health, product quality, and process stability in real time. As fleets of
industrial assets are instrumented with vibration, thermal, acoustic, and current sensors,
the resulting telemetry streams grow far faster than the operations staff available to
interpret them. Anomaly detection --- the task of flagging measurements that deviate from
expected nominal behaviour --- has therefore become a central capability of predictive
maintenance programs, with direct consequences for unplanned downtime, warranty costs, and
worker safety.

Despite substantial research attention, deploying anomaly detection at industrial scale
remains difficult for three interacting reasons. First, sensor fleets are heterogeneous:
different asset classes exhibit markedly different nominal signatures, and a single global
threshold rarely suits all of them. Second, network and compute budgets are constrained at
the edge, so naively streaming raw telemetry to a central server for scoring is often
uneconomical. Third, operating conditions drift over time --- with production schedules,
ambient temperature, and tooling wear all shifting the definition of ``nominal'' --- which
degrades static models within weeks of deployment.
